% Original analytic teaching diagram, 2026-11-12.
% Counting picture behind the capacity formula. Left: codewords (green dots) drawn from the input law.
% Right: the typical output blocks form a set of about 2^{nH(Y)} elements (large blue region); given a codeword, the
% likely outputs form a set of about 2^{nH(Y|X)} elements (small ellipses). The received block (red) lies in the
% ellipse of the sent codeword; decoding fails if it also falls in a competitor's ellipse. Roughly
% 2^{nH(Y)}/2^{nH(Y|X)} = 2^{nI} ellipses fit without much overlap.
\definecolor{VB}{HTML}{7188A3}\definecolor{VBF}{HTML}{E5EAF0}
\definecolor{VG}{HTML}{829C8B}\definecolor{VGF}{HTML}{E8EEE9}
\definecolor{VY}{HTML}{B4A16F}\definecolor{VYF}{HTML}{F1ECDD}
\definecolor{VR}{HTML}{A7787D}\definecolor{VRF}{HTML}{F0E5E6}
\definecolor{VP}{HTML}{9588A5}\definecolor{VPF}{HTML}{ECE8F0}
\definecolor{VI}{HTML}{3B4252}\definecolor{VM}{HTML}{6B7280}\definecolor{VL}{HTML}{4C4D4F}
\begin{tikzpicture}[x=1mm,y=1mm,>=Stealth,font=\figurefont\fontsize{8}{10}\selectfont,
lab/.style={text=VM,font=\figurefont\fontsize{7.5}{9}\selectfont},
cw/.style={circle,fill=VG,inner sep=0pt,minimum size=1.8mm},
fan/.style={draw=VB,line width=.6pt,fill=white},
map/.style={->,draw=VL,line width=.45pt,shorten >=.6mm,shorten <=.6mm}]
% input side
\draw[draw=VL,line width=.6pt] (0,0) ellipse (12 and 17);
\node[lab] at (0,20.5) {输入块};
\node[cw] (c1) at (-4,9) {}; \node[cw] (c2) at (4,3) {}; \node[cw] (c3) at (-3,-5) {}; \node[cw] (c4) at (5,-11) {};
\node[text=VG,anchor=east] at (-5,9) {$x^n(1)$};
\node[text=VG,anchor=east] at (-4,-5) {$x^n(j)$};
% output side
\draw[draw=VB,fill=VBF,line width=.7pt] (62,0) ellipse (30 and 19);
\node[lab] at (62,22.5) {典型输出块，约 $2^{nH(Y)}$ 个};
\draw[fan] (50,8) ellipse (8 and 5);
\draw[fan] (71,7) ellipse (8 and 5);
\draw[fan] (55,-7) ellipse (8 and 5);
\draw[fan] (75,-9) ellipse (8 and 5);
\draw[map] (c1) -- (42.2,8);
\draw[map] (c2) -- (63.2,7);
\draw[map] (c3) -- (47.2,-7);
\draw[map] (c4) -- (67.2,-9);
\fill[VR] (52.5,9.3) circle (1mm);
\node[text=VR,anchor=south] at (52.5,10.2) {$y^n$};
\node[lab,anchor=north] at (62,-20.5) {每个椭圆约 $2^{nH(Y\mid X)}$ 个输出块};
\end{tikzpicture}
